\section{Introduction}\label{sec:intro}
When does a finite game admit a strategy that wins every round? Approaching certainty and attaining certainty are different questions. They also depend on the model of quantum strategy. Table~\ref{tab:landscape} puts the distinction first. The spatial entry is the result of this paper.

\begin{table}[ht]
\centering\small
\begin{tabular}{@{}p{.53\linewidth}ll@{}}
\toprule
Decision problem & Complete class & Source\\
\midrule
Perfect finite-dimensional strategy & $\Sigma^0_1$ & \cite{MIPRE}\\
Perfect commuting strategy & $\Pi^0_1$ & \cite{Slofstra,MNY}\\
Finite-dimensional value equals one & $\Pi^0_2$ & \cite{MNY}\\
Perfect spatial strategy & $\Sigma^0_2$ & This paper\\
\bottomrule
\end{tabular}
\caption{Finite rational Boolean games; computable many-one reductions. A value is a supremum. A perfect strategy attains one.}\label{tab:landscape}
\end{table}

The finite-dimensional and spatial models use tensor products. The latter allows an arbitrary normal state on infinite-dimensional local Hilbert spaces. The commuting model instead allows the two measurement algebras to act on one Hilbert space and commute. Its compact state-space description makes perfect commuting play equivalent to commuting value one. Spatial play has no such compactness in normalized states.

Our upper bound keeps one Schmidt coefficient from disappearing. After that coefficient is scaled to one, the amplitude operator belongs to a norm-compact family of compact operators. Losing probabilities become homogeneous zero equations. These survive a limit even when some squared Schmidt mass escapes. The surviving state can be normalized again.

The hardness mechanism goes in the other direction. A finite Weyl clock commutes with a trace-class density. Each positive density eigenspace is finite-dimensional, so a determinant forces the clock phase to have finite order. The strategy chooses that order. A deterministic computation reads it and checks an infinite sequence of decidable conditions. This is the existential quantifier, and no computable bound on the game can capture its size. The modular-shift toolkit comes from our earlier paper \cite{Entropy}; here four calibrated bridges connect it to the computation.

\Needspace{13\baselineskip}
\begin{theoremA}\label{thm:A}
Deciding whether a finite rational $0/1$ game has a perfect normal spatial strategy is $\Sigma^0_2$-complete. The same holds when a perfect strategy with finite marginal entropy is required. Hardness holds under the promise $\omega_q(G)=1$, for a family in which no instance has a perfect finite-dimensional strategy.
\end{theoremA}

Finite marginal entropy can be required on both sides. For pure strategies the two marginal entropies agree. For a perfect mixed strategy with one finite marginal entropy, a pure component has finite entropy and is still perfect; see Section~\ref{sec:upper}.

\Needspace{10\baselineskip}
\begin{theoremB}\label{thm:B}
For rational Bell functionals, each of the following is $\Sigma^0_2$-complete in the normal spatial model: attainment of a supplied rational global upper bound; attainment of the supremum, with no bound supplied; and attainment of a supplied rational scalar value. In the first problem, validity of the supplied bound can either be promised or required as part of the decision problem.
\end{theoremB}

\Needspace{12\baselineskip}
\begin{theoremC}\label{thm:C}
Uniformly for integers $\kappa\geq1$, the existence of a perfect pure spatial strategy with largest Schmidt probability at least $1/\kappa$ is a $\Pi^0_1$ property. Their union is $\perf_{\qs}\in\Sigma^0_2$. No total computable function $\kappa(G)$ bounds a sufficient stratum for every spatially perfect game. If $\omega_q(G)=1$ and no perfect spatial strategy exists, the largest Schmidt probability tends to zero along every sequence of pure spatial strategies whose winning probabilities tend to one.
\end{theoremC}

\subsection{Prior work and contribution}
The finite-dimensional entry in Table~\ref{tab:landscape} uses perfect completeness in Theorem~12.6 of Ji, Natarajan, Vidick, Wright and Yuen \cite{MIPRE}, not just a gap statement about a supremum. Upper membership follows by a dimension search and real quantifier elimination. Slofstra's embedding theorem \cite{Slofstra} and the account of Mousavi, Nezhadi and Yuen \cite{MNY} give the commuting entry. Their value-one theorem gives the third entry. The commuting upper bound uses the convergent moment hierarchy of Navascu\'es, Pironio and Ac\'in \cite{NPA}. The operator-solution framework for binary constraint system games originates with Cleve and Mittal \cite{CM}; Cleve, Liu and Slofstra gave the commuting-operator extension \cite{CLS}. Slofstra also established nonclosure of finite-dimensional quantum correlations \cite{SlofstraNonclosed}. These are distinct earlier results. Lin's $\mathsf{MIP}^{co}=\mathsf{coRE}$ theorem concerns gap approximation in the commuting model \cite{Lin}; it does not classify spatial attainment.

Infinite-dimensional spatial correlations were constructed by Coladangelo and Stark \cite{CS}. Kim, Paulsen and Schafhauser prove that perfect spatial play collapses to perfect finite-dimensional play for synchronous games \cite[Corollary 3.11]{KPS}. Lupini and coauthors establish the corresponding equivalence for imitation games \cite[Corollary~5.4]{Imitation}. Our games retain nonsynchronous bridge questions. No YES instance is synchronous or an imitation game, since either condition would imply finite-dimensional perfect play.

Man\v{c}inska's flagged magic-square construction supplies the bridge geometry \cite{Mancinska}. The normal bridge, shared-label refinement, graph products and Boolean rational-row tests were developed in \cite[Lemmas~4.1 and 4.2, Theorem~B, Eq.~(32)]{Entropy}. The versions needed here are restated and proved in Section~\ref{sec:compiler}. That section extends the assembly from a particular presentation to four bridges and arbitrary finite lists of endpoint labels. The rooted four-ary computation and the finite-period reader then give the hardness reduction. The compact-sandwich upper bound does not use this compiler.

Fu, Miller and Slofstra study fixed-scenario correlation membership \cite{FMS}. A remark there asserts, without proof, that fixed-scenario membership in $C_q$ and $C_{qs}$ is in RE; Section~\ref{sec:discussion} records a conditional tension with Theorem~A. We do not identify prescribed-table membership with game perfection.

To our knowledge, Theorems~A and B give the first $\Sigma^0_2$-completeness classifications of these spatial attainment problems. This priority statement is provisional. It rests on the cited hierarchy and correlation papers, their nearby results, and a targeted literature search. It is not an exhaustive bibliographic certificate.

\begin{figure}[ht]
\centering\small
\begin{tabular}{@{}c@{}}
Presentation $\xrightarrow{\text{Lemma~\ref{lem:four}}}$ finite game $G_e$\\[4pt]
Perfect normal strategy $\longrightarrow$ normal realization\\[2pt]
$\longrightarrow$ Fock tree $\longrightarrow$ finite clock sector\\[2pt]
$\longrightarrow$ deterministic tableau $\longrightarrow\forall t\ R(e,n,t)$\\[4pt]
Witness $n$ $\longrightarrow$ Fock/Weyl model\\[2pt]
$\longrightarrow$ scaffold, frame, row ancilla $\longrightarrow$ perfect strategy
\end{tabular}
\caption{Dependencies of the hardness proof. Compactness and Bell
attainment have separate upper-bound arguments.}\label{fig:dependencies}
\end{figure}

\subsection{Proof route}
Figure~\ref{fig:dependencies} summarizes the hardness proof. Section~\ref{sec:upper} proves Theorem~C and the finite-entropy upper bound. Section~\ref{sec:compiler} gives the finite game compiler. Sections~\ref{sec:tree}--\ref{sec:reduction} encode a chosen finite witness and all its checks. Section~\ref{sec:value} proves the value-one promise with four explicit boundary defects. Section~\ref{sec:bell} proves Theorem~B. The verification appendix records what the finite programs check.
